\documentclass[11pt]{article}
\usepackage{amsmath,amsthm,amssymb}
\usepackage[margin=1in]{geometry}
\usepackage{hyperref}
\usepackage{booktabs}
\usepackage{enumitem}

\newtheorem{theorem}{Theorem}
\newtheorem{lemma}[theorem]{Lemma}
\newtheorem{proposition}[theorem]{Proposition}
\newtheorem{corollary}[theorem]{Corollary}
\newtheorem{conjecture}[theorem]{Conjecture}
\theoremstyle{definition}
\newtheorem{definition}[theorem]{Definition}
\newtheorem{remark}[theorem]{Remark}
\newtheorem{observation}[theorem]{Observation}

\newcommand{\Hq}{H_q}
\newcommand{\Sn}{S_n}
\newcommand{\Vlam}{V^\lambda}
\newcommand{\Om}{\Omega}
\newcommand{\hatl}{\widehat\lambda}
\newcommand{\Rp}{{R'}}
\newcommand{\SYT}{\mathrm{SYT}}
\newcommand{\elh}{\ell(\hatl)}
\newcommand{\Eplus}{E^+}
\newcommand{\Eminus}{E^-}
\newcommand{\im}{\mathrm{im}}
\newcommand{\pT}{p_T}
\newcommand{\tr}{\mathrm{tr}}

\title{Per-SYT vanishing at the leftmost factor:\\
       \large $p_T(q) = 0$ whenever pair $(\ell, \ell+1)$ is same-column in $T$}
\author{Clio}
\date{17 May 2026}

\begin{document}
\maketitle

\begin{abstract}
For $\lambda = (\hatl, 1^q) \vdash n$ with $\ell = \ell(\lambda) = \elh + q$, let
$\Om^{(\lambda)} = \Rp_{\ell+1} \Rp_{\ell+2} \cdots \Rp_n$ be the column-$1$
chain product in the Iwahori--Hecke algebra $\Hq(\Sn)$, and let
$\pT(q) := \langle v_T, \Om v_T \rangle$ be the diagonal matrix coefficient
of $\Om$ in the Hoefsmit ($b'\!=\!1$) seminormal basis. We prove the
following per-SYT vanishing result: \emph{if the values $\ell$ and $\ell+1$
sit in the same column of $T$ (equivalently, $S_\ell v_T = 0$), then
$\pT(q) = 0$.} The proof is two lines once the leftmost-factor
identity $\Om = S_\ell \cdot X$ is in place: $\Om v_T \in \im(S_\ell) = \Eplus_\ell$
while $v_T \in \Eminus_\ell$, and the seminormal-basis projection onto
$\Eminus_\ell$ extracts the diagonal coefficient cleanly. The result is a
per-SYT refinement of the universal eigenspace statement
$B^{(\lambda)} = \Om V^\lambda \subseteq \Eplus_\ell$ obtained from the same
leftmost-factor identity (May-13 evening-2, prior memo). We verify the
theorem on seven shapes spanning $\tau = 0$ and $\tau \ge 1$ and use it to
account for 26 out of 108 zero diagonals at the smoking-gun shape
$\lambda = (3,3,1,1,1)$, including the historical mystery tableau
$T_0 = ((1,2,5),(3,4,6),(7),(8),(9))$. The remaining 82 zeros at this shape
arise from a finer mechanism left open for future work.
\end{abstract}

\section{Setup}\label{sec:setup}

We adopt the conventions of the trace-vanishing papers~\cite{trace-fwd,converse,pillar1}.

$\Hq(\Sn)$ is the Iwahori--Hecke algebra over $\mathbb{Q}(q)$ with
generators $T_1, \ldots, T_{n-1}$, braid relations
$T_i T_{i+1} T_i = T_{i+1} T_i T_{i+1}$, commutation
$T_i T_j = T_j T_i$ for $|i-j| \ge 2$, and quadratic
$T_i^2 = (q-1) T_i + q$. We set $S_i := T_i + 1$, so the Hecke
quadratic is equivalent to
\begin{equation}\label{eq:hecke-abs}
   S_i^2 \;=\; (q+1) S_i, \qquad T_i S_i \;=\; q\, S_i \;=\; S_i T_i.
\end{equation}
The image $\im(S_i)$ is the $+q$-eigenspace $\Eplus_i := \ker(T_i - q)$,
and the kernel $\ker(S_i)$ is the $-1$-eigenspace $\Eminus_i := \ker(T_i+1)$.
Both are well-defined on any $\Hq$-module, and on the (finite-dimensional)
Specht module $\Vlam$ one has the direct-sum decomposition
\begin{equation}\label{eq:plus-minus-decomp}
   \Vlam \;=\; \Eplus_i \!\mid_{\Vlam} \,\oplus\, \Eminus_i \!\mid_{\Vlam}
\end{equation}
since $S_i$ is diagonalizable on $\Vlam$ with eigenvalues $\{q+1, 0\}$
(this follows immediately from~\eqref{eq:hecke-abs}: every
$x \in \Vlam$ decomposes as $\frac{1}{q+1} S_i x + (x - \frac{1}{q+1} S_i x)$
with the first summand in $\Eplus_i$ and the second in $\Eminus_i$).

For a SYT $T$ of shape $\lambda$, $v_T$ denotes the Hoefsmit seminormal
basis vector, normalised with $b'\!=\!1$ as in~\cite{converse}. The local
action of $S_i$ on $v_T$ depends only on the relative positions of $i$ and
$i+1$ in $T$:
\begin{itemize}[topsep=0.2em,itemsep=0.1em]
\item \textbf{Same row} (SR): $S_i v_T = (q+1) v_T$, so $v_T \in \Eplus_i$;
\item \textbf{Same column} (SC): $S_i v_T = 0$, so $v_T \in \Eminus_i$;
\item \textbf{Diagonal} (D, neither): $S_i$ acts on the
   $2$-dim block $\{v_T, v_{s_i T}\}$ as a non-degenerate
   matrix; both $v_T$ and $v_{s_i T}$ have nonzero components in
   both $\Eplus_i$ and $\Eminus_i$, but the $\Eminus_i$ component is
   contained in the same block $\mathrm{span}(v_T, v_{s_i T})$.
\end{itemize}

For $k \ge 1$, $\Rp_k := S_{k-1} S_{k-2} \cdots S_1$, with $\Rp_1 = 1$.
The chain product on $V_\lambda$ is
\[
   \Om^{(\lambda)} \;:=\; \Rp_{\ell+1}^{(\lambda)} \Rp_{\ell+2}^{(\lambda)} \cdots \Rp_n^{(\lambda)},
\]
where $\ell = \ell(\lambda)$. We write $B^{(\lambda)} = \Om \Vlam$ for the
chain image and $\pT(q) := \langle v_T, \Om v_T \rangle$ for the
$(T,T)$-diagonal coefficient. We use $\langle v_S, v_T \rangle =
\delta_{S,T}$ as the bilinear pairing implicit in ``matrix coefficient''
(equivalently, the coefficient of $v_T$ in the seminormal expansion of
$\Om v_T$).

\section{Statement}\label{sec:statement}

\begin{theorem}[Per-SYT leftmost-factor sign-kill]\label{thm:main}
Let $\lambda \vdash n$ with $\ell = \ell(\lambda)$, and let
$\Om = \Rp_{\ell+1} \Rp_{\ell+2} \cdots \Rp_n$ on $\Vlam$. For any
$T \in \SYT(\lambda)$, if the values $\ell$ and $\ell+1$ sit in the same
column of $T$ (equivalently, the pair $(\ell, \ell+1)$ is SC at $T$, i.e.,
$S_\ell v_T = 0$), then
\[
   \pT(q) \;=\; \langle v_T, \Om v_T \rangle \;=\; 0
   \quad \text{in } \mathbb{Q}(q).
\]
\end{theorem}

The hypothesis is purely a statement about $T$ (a placement of the values
$\ell, \ell+1$); the conclusion is about the chain operator. The theorem
does not require any restriction on $\lambda$ (such as multi-row, $\tau = 0$,
or staircase shape). It is the per-SYT refinement of the universal
containment $B^{(\lambda)} \subseteq \Eplus_\ell|_{\Vlam}$.

\section{Proof}\label{sec:proof}

\subsection*{The leftmost-factor identity}

\begin{lemma}[Leftmost factor]\label{lem:leftmost}
$\Om = S_\ell \cdot X$ for some $X \in \Hq(\Sn)$. Consequently,
$\Om v \in \im(S_\ell) = \Eplus_\ell|_{\Vlam}$ for every $v \in \Vlam$.
\end{lemma}

\begin{proof}
$\Rp_{\ell+1} = S_\ell S_{\ell-1} \cdots S_1 = S_\ell \cdot \Rp_\ell$. So
\[
   \Om \;=\; \Rp_{\ell+1} \Rp_{\ell+2} \cdots \Rp_n
        \;=\; S_\ell \cdot \bigl( \Rp_\ell \,\Rp_{\ell+2} \cdots \Rp_n \bigr)
        \;=\; S_\ell \cdot X
\]
with $X = \Rp_\ell \Rp_{\ell+2} \cdots \Rp_n$. Then
$\Om v = S_\ell (X v) \in \im(S_\ell)$. Since $\im(S_\ell) \subseteq \Eplus_\ell$
on every module (by $T_\ell S_\ell = q S_\ell$, so any $S_\ell w$ satisfies
$T_\ell (S_\ell w) = q (S_\ell w)$), the conclusion follows.\qed
\end{proof}

\subsection*{The projection argument}

\begin{proof}[Proof of Theorem~\ref{thm:main}]
Suppose pair $(\ell, \ell+1)$ is SC at $T$, so $v_T \in \Eminus_\ell|_{\Vlam}$.
By Lemma~\ref{lem:leftmost}, $\Om v_T \in \Eplus_\ell|_{\Vlam}$.

By~\eqref{eq:plus-minus-decomp}, $V_\lambda = \Eplus_\ell \oplus \Eminus_\ell$.
Let $\pi_-: V_\lambda \to \Eminus_\ell$ be the projection along $\Eplus_\ell$.
Then $\pi_-(\Om v_T) = 0$.

Expand $\Om v_T$ in the seminormal basis:
\begin{equation}\label{eq:expansion}
   \Om v_T \;=\; \sum_{T' \in \SYT(\lambda)} c_{T,T'} \, v_{T'},
\end{equation}
with $c_{T,T'} \in \mathbb{Q}(q)$ and $c_{T,T} = \pT(q)$. Apply $\pi_-$:
\begin{equation}\label{eq:proj-equation}
   0 \;=\; \pi_-(\Om v_T) \;=\; \sum_{T'} c_{T,T'} \, \pi_-(v_{T'}).
\end{equation}

Partition $\SYT(\lambda)$ by the type of pair $(\ell, \ell+1)$:
\begin{itemize}[topsep=0.2em,itemsep=0.1em]
\item $T' \in \SYT^{\mathrm{SR}}$: $v_{T'} \in \Eplus_\ell$, so $\pi_-(v_{T'}) = 0$.
\item $T' \in \SYT^{\mathrm{SC}}$: $v_{T'} \in \Eminus_\ell$, so $\pi_-(v_{T'}) = v_{T'}$.
\item $T' \in \SYT^{\mathrm{D}}$: $v_{T'}$ lies in the $2$-dim block
   $W_{T'} := \mathrm{span}(v_{T'}, v_{s_\ell T'})$, with
   $s_\ell T' \in \SYT^{\mathrm{D}}$ also. The block decomposes as
   $W_{T'} = (W_{T'} \cap \Eplus_\ell) \oplus (W_{T'} \cap \Eminus_\ell)$,
   each summand $1$-dimensional. There exist
   $\alpha_{T'}, \beta_{T'} \in \mathbb{Q}(q)$ (depending only on the local
   $T_\ell$-axial distance at $T'$) with
   $\pi_-(v_{T'}) = \alpha_{T'} v_{T'} + \beta_{T'} v_{s_\ell T'}$.
\end{itemize}

Substituting into~\eqref{eq:proj-equation}:
\begin{equation}\label{eq:proj-expanded}
   0 \;=\; \underbrace{\sum_{T' \in \SYT^{\mathrm{SC}}} c_{T,T'}\, v_{T'}}_{(\star)}
   \;+\; \sum_{T' \in \SYT^{\mathrm{D}}} c_{T,T'}\,\bigl(\alpha_{T'} v_{T'} + \beta_{T'} v_{s_\ell T'}\bigr).
\end{equation}

The seminormal basis $\{v_{T'}\}_{T' \in \SYT(\lambda)}$ is linearly
independent. Read off the coefficient of $v_{T''}$ in~\eqref{eq:proj-expanded}
for each $T'' \in \SYT(\lambda)$:
\begin{itemize}[topsep=0.2em,itemsep=0.1em]
\item If $T'' \in \SYT^{\mathrm{SR}}$: $v_{T''}$ does not appear on the
   right (it is not in any $\Eminus_\ell$ block). Coefficient: $0$. No
   constraint on the $c_{T,\cdot}$.
\item If $T'' \in \SYT^{\mathrm{SC}}$: $v_{T''}$ appears only as the
   $T'' = T'$ term of the $(\star)$ sum (it is not in any D-block since
   $s_\ell T''$ is not an SYT). Coefficient: $c_{T,T''}$.
\item If $T'' \in \SYT^{\mathrm{D}}$: $v_{T''}$ appears as the self-term
   of $T' = T''$ in the D-sum (contributing $c_{T,T''} \alpha_{T''}$) and
   as the partner-term of $T' = s_\ell T''$ (contributing
   $c_{T, s_\ell T''} \beta_{s_\ell T''}$). Coefficient:
   $c_{T,T''} \alpha_{T''} + c_{T, s_\ell T''} \beta_{s_\ell T''}$.
\end{itemize}

Setting each coefficient to zero:
\begin{equation}\label{eq:individual-eqns}
   c_{T,T''} \;=\; 0 \quad \text{for every } T'' \in \SYT^{\mathrm{SC}},
\end{equation}
and one linear relation per D-block on $(c_{T,T''}, c_{T, s_\ell T''})$.

By hypothesis $T \in \SYT^{\mathrm{SC}}$. Substituting $T'' = T$ in
\eqref{eq:individual-eqns}:
\[
   \pT(q) \;=\; c_{T,T} \;=\; 0. \qed
\]
\end{proof}

\section{Remarks}\label{sec:remarks}

\subsection{Why the leftmost factor is structurally privileged}

Lemma~\ref{lem:leftmost} extracts $S_\ell$ as the \emph{leftmost} factor
of $\Om$. The same trick cannot be applied to inner factors $S_j$ with
$j > \ell$, because $S_j$ is sandwiched between sub-words of $\Om$ that
contain $S_{j-1}$ and $S_{j+1}$, and braid/commutation relations do not
permit moving $S_j$ to the leftmost position without picking up
correction terms. The symmetric ``rightmost factor'' identity
$\Om = Y \cdot S_1$ — where $Y = \Rp_{\ell+1} \cdots \Rp_{n-1} \cdot
S_{n-1} S_{n-2} \cdots S_2$ — gives \emph{strict} vanishing
$\Om v_T = 0$ when pair $(1, 2)$ is SC at $T$ (descent-$1$): then
$S_1 v_T = 0$, so $\Om v_T = Y \cdot 0 = 0$. This is a trivial
consequence and is independent of Theorem~\ref{thm:main}; the leftmost
analysis is genuinely deeper because it gives $p_T = 0$ from
\emph{coefficient} information (not from $\Om v_T = 0$, which need
not hold under the SC-at-$(\ell, \ell+1)$ hypothesis).

\subsection{Combination with Pillar~1}

The Pillar~1 meta-theorem~\cite{pillar1} establishes
$B^{(\lambda)} \subseteq \Eminus_j|_{\Vlam}$ for $j = 1, 2, \ldots, \tau(\hatl)$
under the standing hypothesis. Run the same projection argument at each
such $j$: if $v_T \in \Eplus_j$ (i.e., pair $(j, j+1)$ is SR at $T$),
then $c_{T, T''}$ vanishes for every $T'' \in \SYT^{\mathrm{SR}\text{ at }j}$.
In particular, taking $T'' = T$ yields:

\begin{corollary}[Pillar-1-side per-SYT vanishing]\label{cor:pillar1-side}
For $\lambda = (\hatl, 1^q)$ under the standing hypothesis of Pillar~1
(\,$\elh \ge 2$, every row $\hatl_r \ge 2$, $q \ge |\hatl| - 2\elh + 2$\,)
and $\tau := \tau(\hatl)$: if there exists $j \in \{1, \ldots, \tau\}$ with
pair $(j, j+1)$ same-row at $T$, then $\pT(q) = 0$.
\end{corollary}

Combined, the two sides yield:

\begin{corollary}[Necessary support of $\pT \ne 0$]\label{cor:nec-support}
At $\lambda = (\hatl, 1^q)$ in the Pillar-1 regime,
$\pT(q) \ne 0$ implies
\begin{itemize}[topsep=0.2em,itemsep=0.1em]
\item pair $(j, j+1)$ is \emph{not} SR at $T$ for any $j = 1, \ldots, \tau$
   (Corollary~\ref{cor:pillar1-side}), and
\item pair $(\ell, \ell+1)$ is \emph{not} SC at $T$
   (Theorem~\ref{thm:main}).
\end{itemize}
\end{corollary}

For descent-$0$ shapes (i.e., $T$ with pair $(1,2)$ SR), the trace
$\tr_{\Vlam}(\Om) = \sum_T \pT(q)$ is concentrated on $T$'s satisfying
the second bullet point. At $\tau = 0$, this is the only ``free''
constraint and Theorem~\ref{thm:main} is the main per-SYT vanishing
mechanism known.

\section{Computational verification}\label{sec:verify}

We verify Theorem~\ref{thm:main} by enumerating $\SYT(\lambda)$ and
computing $\pT(q)$ symbolically in the Hoefsmit ($b'\!=\!1$) basis, using
the framework from~\cite{converse}. For each shape, we count
$N_{\mathrm{SC}}(\ell) := |\{T : \text{pair }(\ell,\ell+1) \text{ SC at } T\}|$,
$N_0^{\mathrm{SC}}(\ell)$ the count of those with $\pT \equiv 0$, and the
total count $N_0$ of zero diagonals across all $T \in \SYT(\lambda)$. The
theorem predicts $N_{\mathrm{SC}}(\ell) = N_0^{\mathrm{SC}}(\ell)$; the
``gap'' $N_0 - N_0^{\mathrm{SC}}(\ell)$ measures vanishings not explained
by Theorem~\ref{thm:main}.

\begin{table}[h]
\centering\renewcommand{\arraystretch}{1.2}
\caption{Verification of Theorem~\ref{thm:main} on seven shapes spanning
$\tau = 0$ and $\tau \ge 1$. The
$N_{\mathrm{SC}}(\ell) = N_0^{\mathrm{SC}}(\ell)$ equality holds in every
case (consistent with the theorem); no counterexamples were found. The
``$N_0$ (total)'' column counts \emph{all} zero diagonals; the gap
$N_0 - N_0^{\mathrm{SC}}(\ell)$ measures vanishings arising from
mechanisms beyond Theorem~\ref{thm:main} (descent-$1$ trivial vanishing,
$\tau$-side Pillar~1 sign-kill, and the open ``interior cancellation''
phenomenon).}
\begin{tabular}{lcccccc}
\toprule
$\lambda$ & $n$ & $\ell$ & $|\SYT(\lambda)|$ &
   $N_{\mathrm{SC}}(\ell)$ & $N_0^{\mathrm{SC}}(\ell)$ &
   $N_0$ (total) \\
\midrule
$(2,2,1)$     & $5$ & $3$ & $5$   & $2$  & $2$  & $4$   \\
$(3,2,1,1)$   & $7$ & $4$ & $35$  & $8$  & $8$  & $30$  \\
$(4,2,1)$     & $7$ & $3$ & $35$  & $6$  & $6$  & $18$  \\
$(2,2,2,1,1)$ & $8$ & $5$ & $28$  & $12$ & $12$ & $28$  \\
$(3,3,1,1,1)$ & $9$ & $5$ & $120$ & $26$ & $26$ & $108$ \\
$(3,2,2,1,1)$ & $9$ & $5$ & $162$ & $45$ & $45$ & $150$ \\
$(4,3,1,1)$   & $9$ & $4$ & $216$ & $30$ & $30$ & $147$ \\
\bottomrule
\end{tabular}
\end{table}

\noindent{}Shapes $(2,2,1)$ through $(3,2,2,1,1)$ and $(2,2,2,1,1)$ were
verified symbolically over $\mathbb{Q}(q)$. Shape $(4,3,1,1)$ was
verified numerically at $q = 11$ and $q = 13$ over $\mathbb{Q}$ — two
generic non-pole values; a nonzero $p_T \in \mathbb{Q}(q)$ has finitely
many roots, so vanishing at both these values is equivalent to symbolic
vanishing modulo the (finite, known) seminormal denominator. Scripts:
\texttt{scratch/2026-05-17-leftmost-factor-verify/}.

\paragraph{Smoking gun.} For $\lambda = (3,3,1,1,1)$, the historically
mysterious zero diagonal at
$T_0 = ((1,2,5),(3,4,6),(7),(8),(9))$ is now \emph{explained}: pair
$(5,6)$ has $5$ at $(0,2)$ and $6$ at $(1,2)$ — same column — so
Theorem~\ref{thm:main} fires with $\ell = 5$, giving $\pT[T_0] = 0$.

\section{Discussion}\label{sec:discussion}

\subsection{What this theorem explains}

At $\lambda = (3,3,1,1,1)$, the agent sweep of 17 May 2026 found 108
zero diagonals out of 120 SYTs. Theorem~\ref{thm:main} accounts for
26 — the ones where pair $(5,6)$ is same-column — and gives a structural
reason: $v_T \in \Eminus_5$ while $\Om v_T \in \Eplus_5$, so the
seminormal coefficient must vanish.

The remaining $108 - 26 = 82$ zeros split further. A subset are
\emph{descent-$1$} SYTs (pair $(1,2)$ SC), for which $\Om v_T = 0$
already (the rightmost-factor argument of Section~\ref{sec:remarks});
these are $f^{(2,2,1,1,1)} = 14$ tableaux. The remaining
$82 - 14 = 68$ are descent-$0$ SYTs with pair $(5,6)$ D — for example
\[
   T \;=\; ((1,2,3),(4,5,8),(6),(7),(9))
\]
where pair $(5,6)$ is D (5 at $(1,1)$, 6 at $(2,0)$). Theorem~\ref{thm:main}
does not fire for these — yet symbolic computation returns $\pT = 0$.
The mechanism here is a higher-order cancellation in the seminormal
expansion of $\Om v_T$, possibly the ``interior sign-kill'' pattern noted
in the May-16 dream-2 entry. Characterising these is an open question.

\subsection{Relation to other techniques}

Theorem~\ref{thm:main} is to the per-SYT diagonal what the leftmost-factor
memo~\cite[\texttt{leftmost-rightmost-factor-technique}]{memory} is to
the chain image: the image $B^{(\lambda)} = \Om \Vlam$ lies in $\Eplus_\ell$,
and the per-SYT specialization extracts the diagonal coefficient from
that universal containment. The proof uses only:
\begin{enumerate}[label=(\roman*),topsep=0.2em,itemsep=0.1em]
\item the Hecke quadratic in the form $S_\ell^2 = (q+1) S_\ell$, which
   gives diagonalizability~\eqref{eq:plus-minus-decomp};
\item the seminormal-basis local action of $S_\ell$ on $v_T$, $v_{s_\ell T}$.
\end{enumerate}
No external machinery (Jucys-Murphy, Markov trace, KL/canonical, BGG,
1423-avoidance, Goertzen-Williamson, etc.) is invoked. This places
Theorem~\ref{thm:main} squarely in the family of one-line algebraic facts
that arise from Hecke absorption — and indeed extends the May-13
leftmost-rightmost-factor toolkit by one notch (per-SYT, not just $B$).

\subsection{What remains open}

\begin{itemize}[topsep=0.2em,itemsep=0.1em]
\item \textbf{Characterising the 82 unexplained zeros at $(3,3,1,1,1)$.}
   These have pair $(\ell, \ell+1)$ D, so $v_T$ has nonzero
   $\Eminus_\ell$ \emph{and} $\Eplus_\ell$ components, and the
   projection-onto-$\Eminus_\ell$ argument no longer forces $c_{T,T}=0$
   directly. A finer mechanism is needed.
\item \textbf{Iterated leftmost factors.} Can the argument be iterated on
   $X = \Rp_\ell \Rp_{\ell+2} \cdots \Rp_n$, exposing a second-leftmost
   $S_{\ell-1}$? The image of $S_\ell S_{\ell-1}$ is a proper subspace of
   $\Eplus_\ell$, but its description is not transparent. A clean
   characterisation would mechanically extend Theorem~\ref{thm:main}.
\item \textbf{D-block constraints.} The D-block linear relations in the
   proof — one per unordered pair $\{T', s_\ell T'\}$ — are nontrivial
   constraints that the present theorem does \emph{not} extract. They
   relate $c_{T,T''}$ and $c_{T, s_\ell T''}$ for D tableaux $T''$. A
   careful analysis of these constraints might yield further vanishing
   in the D regime, possibly explaining a fraction of the 82.
\end{itemize}

\begin{thebibliography}{9}
\bibitem{trace-fwd}
   Clio, \emph{Trace-vanishing forward direction: corollary of multi-row
   meta + right-absorption}, 16 May 2026.
   \texttt{\textasciitilde/projects/proofs/2026-05-16-trace-vanishing-forward.tex}.

\bibitem{converse}
   Clio, \emph{Converse trace-vanishing: per-SYT positivity and
   column-descent-free SYT existence}, 16 May 2026.
   \texttt{\textasciitilde/projects/proofs/2026-05-17-converse-per-syt-positivity.tex}.

\bibitem{pillar1}
   Clio, \emph{Extended sign-kill for general multi-row $\hatl$: the
   meta-theorem at arbitrary $\elh \ge 2$}, 15 May 2026.
   \texttt{\textasciitilde/projects/proofs/2026-05-15-multi-row-sign-kill-meta.tex}.

\bibitem{memory}
   Clio, \emph{Auto-memory note:
   leftmost-rightmost-factor-technique}, 13 May 2026.
   \texttt{\textasciitilde/.claude/projects/-home-clio/memory/leftmost-rightmost-factor-technique.md}.
\end{thebibliography}

\end{document}
